\section{Conventions and normalizations}
\label{app:conventions}

We use
\begin{equation}
 g^{\mu\nu}=\operatorname{diag}(1,-1,-1,-1),
 \qquad D=4-2\eps,
 \qquad \varepsilon^{0123}=+1.
 \label{eq:metric-and-dimension}
\end{equation}
For a vector $a$, $\not\!a=\gamma_\mu a^\mu$, and
\begin{equation}
 i\sigma^{ab}=-\frac12[\not\!a,\not\!b].
 \label{eq:sigma-convention}
\end{equation}
The BMHV chirality matrix is
\begin{equation}
 \gamma_5=\frac{i}{4!}\varepsilon_{\mu\nu\rho\sigma}
 \bar\gamma^\mu\bar\gamma^\nu\bar\gamma^\rho\bar\gamma^\sigma,
 \qquad
 \{\gamma_5,\bar\gamma^\mu\}=0,
 \qquad
 [\gamma_5,\widehat\gamma^\mu]=0.
 \label{eq:bmhv-gamma5}
\end{equation}
All external hadron momenta, light-cone vectors, and polarization vectors are
four dimensional.  Integration momenta are (D) dimensional.

The physical loop measure is $\dd^D\ell/(2\pi)^D$.  The scalar-integral
normalization used by the IBP and dimensional-recurrence formulas is
\begin{equation}
 [\dd^D\ell]_{\rm IBP}=\frac{\dd^D\ell}{i\pi^{D/2}}.
 \label{eq:ibp-measure}
\end{equation}
For an $L$-fold integral with $N_C$ cuts represented by
$2\pi\mathcal C_+$, the conversion is
\begin{equation}
 I_{\rm phys}
 =\left(\frac{i\pi^{D/2}}{(2\pi)^D}\right)^L
 (2\pi)^{N_C}I_{\rm IBP}.
 \label{eq:ibp-physical-conversion}
\end{equation}
This factor changes neither the cut energy orientation nor the ordinary
propagator prescription.

The $\overline{\mathrm{MS}}$ scale factor used in the existing NLO real
calculation is
\begin{equation}
 \left(\frac{\mu^2e^{\gamma_E}}{4\pi}\right)^\eps.
 \label{eq:msbar-scale-factor}
\end{equation}
It is written explicitly when real, virtual, and collinear terms are combined.

The plus distribution is defined by
\begin{equation}
 \int_0^1\dd w\,\mathcal D_m(w)\varphi(w)
 =\int_0^1\dd w\,
 \frac{\log^m(1-w)}{1-w}[\varphi(w)-\varphi(1)]
 \label{eq:plus-test-function}
\end{equation}
for every smooth test function $\varphi$.

\begin{table}[htbp]
\centering
\small
\begin{tabularx}{\textwidth}{@{}L{0.22\textwidth}Y@{}}
\toprule
Symbol & Meaning\\
\midrule
$(P_A,P_B,P_h)$ & Incoming and identified hadron momenta.\\
$(k_a,k_b,k_c)$ & Incoming and identified parton momenta.\\
$(x_a,x_b,z_h)$ & Collinear momentum fractions.\\
$(S,T,U)$ & Hadronic Mandelstam invariants.\\
$(s,t,u)$ & Partonic Mandelstam invariants.\\
$(q,Q^2)$ & Recoil momentum and $Q^2=q^2=s+t+u$.\\
$(v,w,x)$ & Single-inclusive variables with $x=1-w$.\\
$\mathcal C_+^{(n)}$ & Positive-energy cut of propagator power $n$.\\
$\mathcal D_m(w)$ & $[\log^m(1-w)/(1-w)]_+$.\\
$H^{\mathsf P,(n)}$ & Partonic hard coefficient multiplying
$\alpha_s^{n+2}$.\\
\bottomrule
\end{tabularx}
\caption{Principal symbols used in the manuscript.}
\label{tab:symbols}
\end{table}
